%%% -------------------- Titulní strana --------------------
\begin{titlepage}
    \begin{center}
        \large \vspace*{\fill}
        \thispagestyle{empty}

        \LARGE

        { \huge \textbf{Gymnázium Arabská, Praha 6, Arabská 14}}

        {\LARGE Obor programování }

        \vfill
        \includegraphics{logogyarab.png}
        \vspace{15pt}

        \vfill

        {\huge \textbf{Neuronová síť: Digitor}}

        \vfill

        Ondřej Salát

        \vfill

        {\large Květen, 2024}

        \vspace*{\fill}
    \end{center}
\end{titlepage}

%%% -------------------- Prohlášení --------------------
\thispagestyle{empty}
\addtocounter{page}{-1}
\vspace*{\fill}
Prohlašuji, že jsem jediným autorem tohoto projektu, všechny citace jsou řádně označené a všechna
použitá literatura a další zdroje jsou v práci uvedené. Tímto dle zákona 121/2000 Sb. (tzv. Autorský zákon)
ve znění pozdějších předpisů uděluji bezúplatně škole Gymnázium, Praha 6, Arabská 14 oprávnění k výkonu
práva na rozmnožování díla (§ 13) a práva na sdělování díla veřejnosti (§ 18) na dobu časově neomezenou a
bez omezení územního rozsahu.

\vspace{2cm}
V .................. dne .......................

\rightline{Ondřej Salát \hspace{0.5cm} ............................ \hspace{2cm}}

\vspace{2cm}

%%% -------------------- Anotace --------------------
\newpage
\begin{abstract}
    Cílem maturitní práce je vytvořit neuronovou síť a naprogramovat k ní základní učící algoritmus bez použití knihoven pro strojové učení (jako TensorFlow, nebo PyTorch).
    Záměrem je hlubší pochopení a prozkoumání různých implementací neuronových sítí. Neuronová síť Digitor bude natrénována pro rozpoznávání rukou napsaných číslic [0-9].
    Síť dostane na vstup čtvercový obrázek a výstupem bude číslice obsažena v obrázku.
\end{abstract}

%%% -------------------- Obsah --------------------
\tableofcontents